% 张三慧《大学物理学（第4版）》V2，第21章《麦克斯韦方程组和电磁辐射》。
% 主教材扫描PDF物理页235--237，原书印刷页218--220；章末习题21.1--21.17。
% 图21.16--21.18依赖原图，精确TikZ重绘与物理专家复核待办。
\section*{第21章\quad 麦克斯韦方程组和电磁辐射}
\subsection*{章末习题21.1--21.17}

\paragraph{习题 21.1}
利用麦克斯韦方程组证明电荷守恒：若没有电流流入给定的封闭体积，
该体积内的总电荷保持恒定。
可考虑安培--麦克斯韦方程的任意闭合曲面，并令边界环路收缩为一点。

\paragraph{习题 21.2}
对球对称分布的电流场（如点电源向四周均匀发射带电粒子，
或带电球形电容器的均匀径向漏电流），
利用麦克斯韦方程组证明各处磁场 \(\boldsymbol B=\boldsymbol 0\)。
原书图21.16给出径向电流场示意。

\paragraph{习题 21.3}
电子与原子碰撞时获得
\(2.0\times10^{24}\,\mathrm{m/s^2}\) 的减速度。
在与减速度方向夹角 \(45^\circ\)、距离 \(20\,\mathrm{cm}\) 处，
计算电子瞬时辐射横向电场；
再求从碰撞发生至该场信号到达观测点的时间。

\paragraph{习题 21.4}
一束高速电子在X射线管中撞击金属靶；
动能为 \(2\times10^4\,\mathrm{eV}\) 的电子
在 \(5\times10^{-9}\,\mathrm m\) 距离内停止。
取均匀减速近似，求在垂直减速度方向、
距碰撞点 \(0.3\,\mathrm m\) 处辐射电场的大小。

\paragraph{习题 21.5}
无线上天线（一段直导线）内电子作简谐振动，
速度 \(v=v_0\cos\omega t\)，其中
\(v_0=8.0\times10^{-3}\,\mathrm{m/s}\)，
\(\omega=6.0\times10^6\,\mathrm{rad/s}\)。
(1) 求电子最大加速度；
(2) 求沿垂直天线方向、距天线 \(1.0\,\mathrm{km}\) 处
单电子辐射横向电场的最大值，
及从最大加速度时刻至信号抵达的延迟。

\paragraph{习题 21.6}
在范德格拉夫加速器中，质子获得
\(1.1\times10^{14}\,\mathrm{m/s^2}\) 的加速度。
(1) 在与加速度夹角 \(45^\circ\)、距质子 \(0.50\,\mathrm m\)
处求横向辐射电场与磁场的大小；
(2) 画出加速度、传播方向、电场、磁场的方向关系。

\paragraph{习题 21.7}
直径 \(0.26\,\mathrm{cm}\) 的铜导线内部电场
\(3.9\times10^{-3}\,\mathrm{V/m}\)，
流过恒定电流 \(12.0\,\mathrm A\)。
(1) 求单个自由电子由该场引起的加速度；
(2) 求该电子在垂直导线方向、距其 \(4.0\,\mathrm m\)
处辐射横向电场和磁场；
(3) 假设长 \(5.0\,\mathrm{cm}\) 的导线段内所有自由电子
同时产生电磁辐射，估算 \(4.0\,\mathrm m\) 处总横向场；
(4) 求同段 \(12.0\,\mathrm A\) 恒定电流产生的磁场并讨论
(3) 的结果为何不能直接观测。
\textit{题设的相干叠加是假想条件；须讨论自由电子散射、相位与稳恒电流。}

\paragraph{习题 21.8}
太阳光在地球大气层顶的辐射强度
\(I=1.38\times10^3\,\mathrm{W/m^2}\)。
将太阳光近似为简谐电磁波，
求其中电场强度和磁感应强度的方均根值。

\paragraph{习题 21.9}
用于打孔的激光束截面直径 \(60\,\mathrm{\mu m}\)，
总功率 \(300\,\mathrm{kW}\)。
求光束的坡印亭矢量时间平均大小，
以及电场、磁场的振幅。

\paragraph{习题 21.10}
波长 \(514.5\,\mathrm{nm}\) 的氩离子激光器，
以 \(3.8\,\mathrm{kW}\) 向月球发射光束。
光束全发散角为 \(0.880\,\mathrm{\mu rad}\)，
地月距离取 \(3.82\times10^5\,\mathrm{km}\)。
求 (1) 月面光斑圆面积的半径；
(2) 到达月面时的平均辐射强度。
\textit{按理想均匀远场光斑估算，不计大气与衍射结构细节。}

\paragraph{习题 21.11}
圆柱形导体长 \(l\)、半径 \(a\)、电阻率 \(\rho\)，
通电流 \(I\)，表面不带电（图21.17）。
证明 (1) 导体表面处的坡印亭矢量
处处与表面垂直且指向导体内部；
(2) 通过导体侧面的\emph{向内}能流总功率
为 \(I^2R\)，其中 \(R=\rho l/(\pi a^2)\)。
\textit{使用外法线时，坡印亭通量积分应为 \(-I^2R\)。}

\paragraph{习题 21.12}
同轴电缆内芯与外金属筒半径为
\(r_1,r_2\)，由电源电动势 \(\mathcal E\)
向电阻 \(R\) 送电（图21.18），
电缆中间为真空。
(1) 求 \(r_1<r<r_2\) 区域的
\(\boldsymbol E,\boldsymbol B,\boldsymbol S\)
及能流方向；
(2) 在横截面环形区域积分证明传输功率
\(\mathcal E^2/R\)，正是负载耗散的功率。
\textit{取稳定理想同轴几何并忽略线损和边缘场。}

\paragraph{习题 21.13}
平面简谐电磁波的波长 \(3.0\,\mathrm{cm}\)，
电场振幅 \(30\,\mathrm{V/m}\)。
求 (1) 频率；(2) 磁场振幅；
(3) 面积 \(0.5\,\mathrm{m^2}\)、
垂直传播方向的完全吸收表面所受平均辐射压力，
并给出相应总作用力。

\paragraph{习题 21.14}
太阳直射海滩的强度为
\(1.1\times10^3\,\mathrm{W/m^2}\)。
迎光面积 \(0.5\,\mathrm{m^2}\)，
皮肤反射率 \(50\%\)。
求受到的平均辐射压力及总力，
假设镜面正入射、其余光被吸收。

\paragraph{习题 21.15}
强激光束压缩高电子密度等离子体，
假设能将入射激光全部反射。
脉冲峰值功率 \(1.5\times10^9\,\mathrm W\)，
照射面积 \(1.3\,\mathrm{mm^2}\)。
求等离子体承受的峰值辐射压强。

\paragraph{习题 21.16}
宇航员离太空舱 \(10\,\mathrm m\)，持有向后发射的
\(10\,\mathrm{kW}\) 激光枪；
连同所携物品总质量 \(100\,\mathrm{kg}\)。
若光束始终朝远离太空舱方向射出，
求光反冲使其从静止回到舱边的时间。
忽略其他力和激光器的反作用质量变化。

\paragraph{习题 21.17}
在绕太阳的圆轨道上有均匀球形尘埃，
密度 \(1.0\,\mathrm{g/cm^3}\)。
太阳表面辐射强度
\(6.9\times10^7\,\mathrm{W/m^2}\)。
若尘埃完全吸收阳光，求其半径达到多少时，
向外的辐射压力恰好平衡太阳向内的万有引力；
此时其轨道运动会怎样变化？
取太阳质量与半径的标准近似值并注明。
